%% kapa: exam=1 round=37 year=2026 period=1 form=A subject=real_estate subject_ko=부동산학원론 seq=15 number=95
%% source: Q-Net 감정평가사 기출문제 (cst00309) artlSeq=5258924 file=raw/1008/5258924/2026년 제37회 감정평가사 1차 1교시.pdf
%% conversion: kordoc 4.24.51 -> tools/gen1.py (auto); review=proofread-sonnet; flags=none
\begin{question}{95}
부동산시장의 특성으로 옳지 않은 것은?
\begin{choices}
  \item 지리적 위치의 고정성으로 인해 부동산시장은 지역에 따라 여러 개의 부분시장으로 세분된다.
  \item 시장상황의 변화에 따른 수급조절이 어렵기 때문에 부동산시장에서는 단기적으로 가격의 왜곡이 발생할 가능성이 높다.
  \item 부동산에는 법적 제한이 많기 때문에 부동산시장이 왜곡되고 시장의 조절기능이 저하된다.
  \item 부동산은 대체가능성이 낮은 재화이기에 부동산시장에서는 공매(short selling)가 빈번하게 발생한다.
  \item 부동산이 고가의 상품이기 때문에 부동산시장은 자금의 유용성에 큰 영향을 받는다.
\end{choices}
\end{question}
\answer{4}
